% [原文 p.v]
\chapter*{引言}
\addcontentsline{toc}{chapter}{引言}
\markboth{引言}{引言}

本书基于 1964 年秋在哈佛大学所授的课程。该课程旨在给出向量丛与 K 理论的一个自足的说明，只以拓扑空间论与线性代数的基础为前提。本书处理方式的一个特点是，不以通常的同调论或上同调论为中心。事实上，有理（系数）上同调正是借助 K 理论来定义的。

在 K 理论中，讨论一直延伸到 Hopf 不变量问题的解决，也涉及 $J$ 同态的初步内容。除讲义之外，本书还把 1964 年以后发表的两篇论文作为附录 B、C 重新收录。其中一篇讨论幂算子，适合作第 3 章题材的补充：与第 3 章中处理算子的 Grothendieck 手法相比，它并不精致，但提供了另一条基本途径，把幂算子与滤过联系起来。实际上，讲义处理的是紧空间而不是胞腔复形，所以在讲义中不解释由切片给出的滤过。另一篇论文给出看待 $KR$ 理论的新视角，弥补了讲义中明显的不足。

% [Tutor Note: 日译本此处两处都只写「作用素」，原书英文为 power operations，指的是附录 B 的幂算子（$\psi^{k}$ 一类），与第 3 章讨论的一般算子不同；日译本在此省略了「冪」字，译文按原书英文补足。]
% CHUNK fm DONE
